在\hyperref[sec:09-Convex-Optimization/07-Applications/03]{``最大间隔与支持向量机''}的结尾，对偶问题暴露出一个容易被放过的观察：训练数据从不单独出现，只以内积 \(\ip{x_i}{x_j}\) 的形式出现。把内积换成某个更一般的函数 \(k(x_i,x_j)\)，线性分离面就悄悄搬到了另一个空间。这个观察能走多远，能解释多少已经学过的方法，又要付出什么代价？

先看一个具体的数字。在 \(x\in\R^2\) 上做 Ridge 回归，用二次多项式核 \(k(x,z)=(1+\ip{x}{z})^2\) 代替线性内积。取 \(x=(1,2)\)，\(z=(3,4)\)：

\[
k(x,z)=(1+1\cdot3+2\cdot4)^2=(1+3+8)^2=144.
\]

这个值等于把 \(x\) 与 \(z\) 映射到 \(\phi(x)=(1,\sqrt2 x_1,\sqrt2 x_2,x_1^2,x_2^2,\sqrt2 x_1x_2)\) 之后做六维内积，但计算只用了两次乘法与一次平方。特征映射存在，却从未被构造出来。核技巧的价值就在这里：算法只要看得见内积，就能在极高维甚至无穷维的特征空间里工作，而不必把那个空间造出来。

这一单元沿着这条线索推进，五节的关系可以先看一眼全貌。

\begin{center}
\begin{tikzpicture}[>=stealth,font=\footnotesize,
  bx/.style={draw,rounded corners,align=center,
             minimum width=2.9cm,minimum height=1.0cm}]

\node[bx] (A) at (0,1.8) {SVM 对偶\\数据只以内积出现};
\node[bx] (B) at (4.5,1.8) {核技巧\\内积换成 \(k\)};
\node[bx] (C) at (9.0,1.8) {RKHS 与表示定理\\无穷维塌缩成 \(n\) 维};
\node[bx] (D) at (6.6,-0.4) {核岭回归\\闭式解与 \(O(n^{3})\)};
\node[bx] (E) at (11.4,-0.4) {Gaussian process\\均值与方差};

\draw[->] (A) -- (B);
\draw[->] (B) -- (C);
\draw[->] (C) -- (D);
\draw[->] (C) -- (E);

\node[align=center,black!60] at (2.25,0.9) {半正定是\\通行证};
\draw[->,black!45] (2.25,1.25) -- (2.25,1.55);

\end{tikzpicture}
\end{center}

\emph{同一个 \(k\)，先当内积用，再当范数用，最后当协方差用。}

起点是线性模型的表达力边界：有些数据在任何直线、任何平面下都分不开，升到高维后却变得线性可分，于是``升到多高、怎样升''成了矛盾所在。出路是不显式构造高维特征，只通过内积访问它。随后必须回答一个理论问题：随手写出的相似度函数，什么时候真的是某个空间里的内积？半正定性给出判据，再生核希尔伯特空间给出构造，表示定理则把无穷维的函数优化塌缩成 \(n\) 维的系数问题。最后落到计算：核岭回归给出闭式解，Gram 矩阵的存储与求逆代价划定规模边界；而把同一个 \(K\) 读成协方差，预测就从一个点变成一个分布。

\begin{enumerate}
\tightlist
\item
  把数据升到高维为什么有用：XOR 这类数据在低维线性不可分，显式特征映射让它可分，但维度爆炸使显式构造不可行，逼出``只算内积''的出路。
\item
  核技巧：只通过内积看世界：定义核函数，盘点 SVM 对偶、Ridge 回归的 Gram 形式与 PCA 中只依赖内积的环节，给出常见核，并说明自造相似度不满足半正定时会坏在哪里。
\item
  再生核希尔伯特空间：把函数当成点：从 \(k(x,\cdot)\) 张成的空间出发构造内积与完备化，说明求值泛函连续意味着什么，并用表示定理把无穷维优化塌缩成 \(n\) 维问题。
\item
  核岭回归与核方法的代价：用表示定理推出闭式解，与 Ridge 的原始形式对照，并算清 \(O(n^{2})\) 存储、\(O(n^{3})\) 求解与预测时保留全部训练数据这三笔账。
\item
  Gaussian process 把核变成函数先验：同一个核不再只定义相似度与范数，还定义函数值的联合协方差；Gaussian 条件化给出均值与方差，而方差只反映输入位置的疏密，不反映模型设定是否正确。
\end{enumerate}

\subsection{怎么读这一单元}

只想快速拿起这套方法的话，读第一节的 XOR、第二节的替换规则、第四节的两笔规模账、第五节的均值与方差即可，这条线回答``为什么有用、怎样计算、何时不可信''。想弄清方法为何成立，再补上半正定判据、再生性质与表示定理的那三行正交分解。要实现或建模时，最后检查核的选择、正则强度、数值条件、样本规模与分布外外推，这些边界比记住某个闭式解更有用。

阅读时可以始终追问四个问题：\(\phi\) 到底是谁、是否真的存在？每一步计算是否只用到内积？范数惩罚究竟在惩罚什么？样本量翻倍时，哪一步先撑不住？图中那条从左到右的链条，正是在这四个问题上反复接受检验的结果。

学完之后建议回到\hyperref[ch:118]{``统计学习理论：训练误差对真风险有多乐观''}，看由 \(\norm{f}_{\mathcal{H}_k}\) 控制的函数类在泛化界中扮演什么角色；那一章会说明，这里反复出现的那个范数为何值得付出 \(O(n^{3})\) 的代价。
